\chapter{Metodolog\'ia}

Este cap\'itulo define c\'omo se estudia el sistema y con qu\'e criterio se juzgan sus resultados. No describe el funcionamiento interno del flujo, que se
desarrolla en el cap\'itulo siguiente, sino el marco con el que se analiza su comportamiento de una forma defendible. La idea central es valorar la
viabilidad de una arquitectura basada en LLM que prioriza la trazabilidad y la explicabilidad frente a una respuesta generativa directa y opaca.

\section{Enfoque metodol\'ogico}

El trabajo se plantea como una evaluaci\'on emp\'irica de un sistema orientado al an\'alisis financiero hist\'orico a partir de consultas en lenguaje natural.
El inter\'es principal no es medir la m\'axima capacidad posible de un modelo aislado, sino estudiar si una arquitectura m\'as controlada, compuesta por
fases separadas y artefactos observables, puede producir respuestas \'utiles, comprensibles y prudentes para el usuario final.

Desde este punto de vista, la metodolog\'ia no se centra solo en si el sistema ``responde'', sino en dos cuestiones complementarias. La primera es si logra
completar una ejecuci\'on de extremo a extremo y entregar una salida final utilizable. La segunda es si esa respuesta final mantiene una calidad
suficiente en relaci\'on con lo que ped\'ia la consulta. Esta doble lectura obliga a distinguir entre completar t\'ecnicamente el flujo y obtener un
resultado satisfactorio desde la perspectiva del usuario.

\section{Pregunta de investigaci\'on y objetivo evaluable}

La pregunta que orienta este cap\'itulo puede formularse del siguiente modo: \textit{evaluar la viabilidad de una arquitectura LLM trazable y explicable
para resolver consultas de an\'alisis financiero hist\'orico en lenguaje natural.}

Esta formulaci\'on es coherente con el objetivo general del trabajo, que no busca maximizar la capacidad bruta de un modelo poco interpretable, sino
comprobar si es posible construir un sistema m\'as controlado, capaz de dejar evidencias del proceso seguido y de acotar mejor d\'onde y por qu\'e se
producen los errores.

En consecuencia, el objetivo evaluable no consiste en demostrar que el sistema puede responder a cualquier consulta financiera, ni en compararlo con un
analista experto, ni en presentarlo como herramienta de predicci\'on o asesoramiento. Lo que se quiere comprobar es si el flujo propuesto resulta
viable como sistema de apoyo para an\'alisis hist\'orico: esto es, si puede resolver un conjunto representativo de consultas, mantener un comportamiento
trazable y producir respuestas finales suficientemente adecuadas para el uso planteado.

\section{Unidad de an\'alisis}

La unidad de an\'alisis principal es una consulta completa ejecutada de extremo a extremo. Esta decisi\'on implica que el caso evaluado no es una etapa aislada
del workflow ni un artefacto intermedio concreto, sino el recorrido completo que va desde la petici\'on inicial hasta la respuesta final.

Esta elecci\'on responde a una raz\'on simple. El sistema est\'a pensado para que el usuario formule una consulta y reciba una respuesta final. Por tanto, no
tiene sentido considerar como caso satisfactorio una ejecuci\'on que no llega a producir una salida utilizable para el usuario, aunque algunas fases internas
hayan funcionado correctamente. Del mismo modo, tampoco basta con llegar al final del flujo si la respuesta obtenida no mantiene una calidad suficiente en
relaci\'on con la consulta planteada.

\section{Dise\~no experimental}

La evaluaci\'on se apoya en la bater\'ia de 50 consultas descrita en el cap\'itulo anterior. Esa bater\'ia se organiza en tres niveles de complejidad,
A, B y C, y se distribuye de forma deliberadamente exigente en 10 casos de nivel A, 20 de nivel B y 20 de nivel C. Esta estructura permite observar no
solo el comportamiento del sistema en consultas relativamente directas, sino tambi\'en su respuesta cuando aumentan las exigencias anal\'iticas, de formato o
de prudencia interpretativa.

Cada caso se ejecuta una \'unica vez y todos los ejemplos se lanzan con la misma configuraci\'on del sistema y del modelo. En particular, la evaluaci\'on
se realiza con el modelo \texttt{openai/gpt-oss-20b} \cite{openai2025gptossmodelcard}, manteniendo constante la configuraci\'on general para toda la bater\'ia. De 
este modo, las diferencias observadas entre casos pueden interpretarse en funci\'on de la dificultad de la consulta y del comportamiento del flujo, y no como 
consecuencia de cambiar de modelo o de criterio de ejecuci\'on entre ejemplos.

\section{Protocolo general de ejecuci\'on}

El protocolo de evaluaci\'on sigue una misma l\'ogica para todos los casos. A partir de una consulta en lenguaje natural, el sistema intenta completar el
recorrido de extremo a extremo: resoluci\'on de datos, descarga, an\'alisis, generaci\'on y validaci\'on del c\'odigo, ejecuci\'on e interpretaci\'on final.
Si el flujo consigue terminar y entregar una respuesta final utilizable, el caso pasa a la fase de evaluaci\'on cualitativa. Si no lo consigue, el caso se
registra como no completado y se analiza aparte seg\'un la familia de error observada.

Aunque el tiempo de ejecuci\'on de cada caso se conserva como informaci\'on complementaria, su interpretaci\'on principal se reserva para el cap\'itulo de resultados.

\section{Criterio general de \'exito}

Un caso no se considera exitoso solo por llegar a una respuesta final. Para que el resultado se considere satisfactorio dentro del trabajo deben cumplirse dos
condiciones. La primera es que el flujo complete la ejecuci\'on y produzca una salida final utilizable. La segunda es que esa respuesta alcance una calidad
suficiente en la revisi\'on manual posterior.

En este sentido, se adopta como criterio de referencia que un resultado satisfactorio es aquel que obtiene al menos 8 puntos sobre 10 en la r\'ubrica de evaluaci\'on. 
Alcanzar ese umbral significa que la respuesta final presenta una calidad suficientemente alta como para considerarla adecuada, clara y prudente dentro del alcance 
del sistema. Las respuestas que completan el flujo pero no llegan a ese nivel no se consideran fallos t\'ecnicos del mismo tipo
que los casos no completados, pero s\'i se interpretan como resultados que todav\'ia requieren mejora para alcanzar una calidad suficiente.

\section{Evaluaci\'on de casos no completados}

Los casos que no consiguen producir una respuesta final utilizable no se eval\'uan con la r\'ubrica cualitativa. En su lugar, se analizan por separado seg\'un la 
familia de fallo observada. Esta clasificaci\'on no pretende fijar una taxonom\'ia cerrada desde el principio, sino ofrecer un marco inicial ordenado para estudiar
d\'onde se concentran los bloqueos del sistema. Como familias generales de partida se consideran las siguientes:

\begin{itemize}
    \item fallo en resoluci\'on o descarga de datos;
    \item fallo en generaci\'on o validaci\'on del an\'alisis;
    \item fallo de ejecuci\'on;
    \item fallo en la obtenci\'on de una salida final utilizable.
\end{itemize}

La raz\'on para trabajar con familias y no solo con errores aislados es que dos casos distintos pueden fallar por mensajes concretos diferentes y, sin embargo,
pertenecer al mismo tipo de problema metodol\'ogicamente relevante. Esta agrupaci\'on permite analizar tendencias generales sin perder de vista que, al
revisar los resultados reales de la bater\'ia, podr\'an aparecer matices o subtipos adicionales. Por tanto, las familias de error se fijan aqu\'i como un
marco inicial abierto a refinarse si los resultados muestran situaciones nuevas o diferencias que merezcan una separaci\'on m\'as precisa.

\section{Evaluaci\'on de casos completados}

Solo los casos que llegan a una respuesta final pasan a la evaluaci\'on cualitativa. En ellos, el inter\'es principal no es inspeccionar artefactos
internos que el usuario no ve, sino juzgar la calidad de la salida final que el sistema entrega.

La r\'ubrica se apoya en cinco criterios, todos ellos ponderados por igual:

\begin{itemize}
    \item adecuaci\'on a la consulta;
    \item cobertura anal\'itica;
    \item coherencia y solidez aparente;
    \item claridad y utilidad comunicativa;
    \item prudencia y tratamiento de limitaciones.
\end{itemize}

El primer criterio valora si la respuesta responde realmente a lo que ped\'ia el usuario, incluyendo activos, periodo, comparaci\'on, tono o formato cuando
esos elementos se exigen de manera expl\'icita. El segundo comprueba si la respuesta cubre los componentes anal\'iticos esenciales de la consulta, sin
premiar la longitud por s\'i misma. El tercero se centra en la coherencia interna del resultado y en si las cifras, relaciones y conclusiones aparecen
integradas de forma consistente. El cuarto estudia si la salida es comprensible y \'util para un lector no especializado. El quinto revisa si la respuesta
mantiene la prudencia exigible a un sistema de an\'alisis hist\'orico y si trata correctamente sus limitaciones.

Para que esta evaluaci\'on resulte interpretable, cada criterio debe entenderse de forma operativa. No se trata de asignar una impresi\'on general difusa, sino
de comprobar qu\'e aparece realmente en la respuesta final y c\'omo se ajusta a la consulta planteada.

\begin{itemize}
    \item \textbf{Adecuaci\'on a la consulta.} Eval\'ua si la respuesta contesta a lo que el usuario pidi\'o. Aqu\'i se revisa si aparecen los activos
    correctos, el periodo o comparaci\'on solicitados y, cuando procede, el formato pedido por el usuario. Se asigna un 0 cuando la respuesta falla en
    elementos esenciales de la consulta; un 1 cuando responde solo de forma parcial o altera alguna parte importante; y un 2 cuando responde de manera
    clara y ajustada a los requisitos principales del caso.

    \item \textbf{Cobertura anal\'itica.} Eval\'ua si la respuesta incluye los elementos anal\'iticos que hac\'ian falta para contestar bien a la consulta.
    No premia decir muchas cosas, sino cubrir lo necesario. Se asigna un 0 cuando faltan componentes esenciales del an\'alisis; un 1 cuando aparece una
    parte relevante pero se omite alg\'un elemento importante; y un 2 cuando la respuesta recoge de forma suficiente las m\'etricas, comparaciones o bloques
    que el caso exig\'ia.

    \item \textbf{Coherencia y solidez aparente.} Eval\'ua si la respuesta mantiene coherencia interna y si sus cifras, relaciones y conclusiones
    parecen bien integradas. No equivale a una auditor\'ia financiera experta, pero s\'i exige consistencia visible. Se asigna un 0 cuando hay
    contradicciones claras, cifras mal integradas o conclusiones incoherentes; un 1 cuando la respuesta es en general razonable, pero presenta alguna
    imprecisi\'on o una relaci\'on d\'ebil entre datos y conclusi\'on; y un 2 cuando el resultado aparece consistente y bien construido desde el punto de
    vista del lector.

    \item \textbf{Claridad y utilidad comunicativa.} Eval\'ua si la respuesta puede entenderse con facilidad y si resulta \'util para un lector no
    especializado. Se asigna un 0 cuando la respuesta es confusa, desordenada o poco aprovechable; un 1 cuando se entiende, pero con carencias de orden o
    explicaci\'on; y un 2 cuando la salida es clara, legible y transmite una idea \'util para interpretar el resultado.

    \item \textbf{Prudencia y tratamiento de limitaciones.} Eval\'ua si la respuesta respeta el alcance del sistema, evita recomendaciones o
    predicciones impropias y trata correctamente las limitaciones del an\'alisis hist\'orico. Se asigna un 0 cuando incurre en afirmaciones que exceden ese
    alcance; un 1 cuando mantiene cierta prudencia, pero trata las limitaciones de forma insuficiente; y un 2 cuando conserva un tono prudente, reconoce
    los l\'imites relevantes y no presenta el resultado como asesoramiento financiero.
\end{itemize}

\section{Escala de puntuaci\'on y r\'ubrica}

Cada uno de los cinco criterios anteriores se punt\'ua con una escala discreta de 0, 1 o 2. Esta elecci\'on se adopta por razones de interpretabilidad y
consistencia. Una escala corta facilita distinguir entre incumplimiento claro, cumplimiento parcial y cumplimiento adecuado, y evita una falsa precisi\'on que
ser\'ia dif\'icil de sostener en una evaluaci\'on humana de este tipo.

En esta escala, el valor 0 corresponde a un incumplimiento claro del criterio; el valor 1 indica un cumplimiento parcial o insuficiente; y el valor 2
representa un cumplimiento adecuado. Esta gradaci\'on permite objetivar en la medida de lo posible una revisi\'on que sigue siendo humana, ya que obliga a
distinguir entre respuestas claramente incorrectas, respuestas aceptables pero todav\'ia mejorables y respuestas que satisfacen bien el criterio evaluado.

La puntuaci\'on total m\'axima es de 10 puntos. Todos los criterios pesan lo mismo, ya que el objetivo no es imponer una ponderaci\'on artificial entre
dimensiones, sino obtener una lectura equilibrada de la calidad final de la respuesta. En este marco, los resultados se interpretan de la siguiente forma:

\begin{itemize}
    \item de 8 a 10 puntos: resultado satisfactorio;
    \item de 6 a 7 puntos: caso completado con calidad insuficiente;
    \item de 0 a 5 puntos: caso completado con calidad baja.
\end{itemize}

De esta manera, la metodolog\'ia separa con claridad tres situaciones distintas: casos no completados, casos completados pero todav\'ia insuficientes
y casos completados con una calidad final satisfactoria.

\section{Revisi\'on humana}

La aplicaci\'on de la r\'ubrica se realiza mediante revisi\'on humana llevada a cabo por el autor del trabajo. Esta decisi\'on se justifica porque varios de
los criterios evaluados ---como la adecuaci\'on a la consulta, la claridad de la respuesta o la prudencia interpretativa--- requieren una lectura contextual
que no conviene reducir a una \'unica comprobaci\'on autom\'atica.

Al mismo tiempo, esta revisi\'on no debe interpretarse como una auditor\'ia financiera experta ni como una certificaci\'on profesional del contenido
econ\'omico. Su funci\'on es valorar el comportamiento del sistema dentro del marco del trabajo y determinar si la respuesta final resulta suficientemente
adecuada para el objetivo planteado. Precisamente por ello, la r\'ubrica se formula con criterios expl\'icitos y una escala simple, buscando que el juicio
manual sea lo m\'as interpretable posible.

\section{Amenazas a la validez}

La primera limitaci\'on metodol\'ogica es que la revisi\'on cualitativa depende del juicio del autor del trabajo. Aunque la r\'ubrica reduce parte de esa
subjetividad al definir criterios expl\'icitos y una escala cerrada, el resultado sigue descansando en una evaluaci\'on humana no experta en finanzas.
Por tanto, las puntuaciones deben leerse como una valoraci\'on aplicada del comportamiento del sistema y no como una auditor\'ia financiera externa.

La segunda limitaci\'on es la dependencia de una \'unica configuraci\'on del sistema y de un \'unico modelo, \texttt{openai/gpt-oss-20b}
\cite{openai2025gptossmodelcard}. Esta decisi\'on favorece la comparabilidad interna de la bater\'ia, pero restringe el alcance de las conclusiones: los
resultados obtenidos no deben generalizarse autom\'aticamente a otras configuraciones, otros proveedores o modelos con distinto perfil de capacidad.

La tercera limitaci\'on procede del propio material de evaluaci\'on. Aunque la bater\'ia de 50 consultas cubre diferentes niveles de dificultad y varios tipos
de instrumento, sigue siendo un conjunto finito y deliberadamente acotado. En consecuencia, la metodolog\'ia permite estudiar la viabilidad del sistema en ese
marco concreto, pero no demostrar un comportamiento universal frente a cualquier consulta financiera posible.

Por \'ultimo, el alcance metodol\'ogico de la evaluaci\'on se limita al an\'alisis hist\'orico. Los resultados no deben extrapolarse a tareas de predicci\'on, 
asesoramiento financiero o interpretaci\'on de factores externos. Las implicaciones generales de este l\'imite se desarrollan con mayor amplitud en el cap\'itulo de limitaciones.
